\chapter{Coherence as Nature's Justice}

\epigraph{``Each part does its own work, and the whole is harmonious.''}{\textcite{plato1997complete},
\emph{Republic} 433a}

\section{Justice in the Soul and City}

In the \emph{Republic}, Socrates discovers justice at the moment he ceases to
look for it as a thing and begins to see it as a relation. Justice is ``each
part doing its own work'' (\emph{Republic} 433a). In the city, justice is the
craftsman practising his craft, the soldier soldiering, the philosopher ruling
--- each fulfilling the function for which their nature fits them. In the soul,
justice is reason governing, spirit supporting, appetite obeying --- each
faculty operating according to its proper excellence.

Justice, for Plato, is not an external standard imposed on the city or the soul.
It is the \emph{internal coherence} of a well-ordered whole, where each part
acts according to its nature and the whole functions as it should.

\section{Justice in Nature}

We extend Plato's discovery: nature, too, has its justice, and it is nothing
other than coherence.

An organism is ``just'' when its parts function together to sustain its life ---
when the heart beats, the lungs breathe, the cells metabolise, each doing its
own work. An ecosystem is ``just'' when its species occupy their niches without
one dominating to the point of collapse. A physical system is ``just'' when its
forces balance such that the system persists.

This is the \emph{coherence} that the \physis\ engine tracks. The three-valued
scale --- +1 (Success), 0 (Inert), -1 (Failure) --- is a measure of the extent
to which an action is ``doing its own work'' in the natural economy. An action
that produces a real, detected effect is coherent with the order of nature. An
action that produces no effect is inert --- it neither helps nor harms but also
does nothing. An action that violates its own declared intent or produces a
harmful effect is a failure, a violation of natural justice.

\section{The Coherence of the Practitioner}

The practitioner is ``just'' in the Platonic sense when their actions are
coherent with their declared intent and with the natural order in which they
act. This has two dimensions.

\subsection{Internal Coherence: The Act Matches the Intent}

The first dimension is the match between what the practitioner declares they
will do and what they actually do. This is the most basic form of coherence:
the practitioner who declares a goal and acts toward it scores +1; the
practitioner who declares a goal and does nothing scores 0; the practitioner
who declares a goal and acts against it scores -1.

The \praxis\ platform registers every declared goal and every registered action
as separate entries in the coherence graph. The engine compares them. The
verdict is not a moral judgment but a \emph{structural} one: does the action
satisfy the conditions set by the intent?

\subsection{External Coherence: The Act Produces Its Intended Effect}

The second dimension is the match between the action and its consequence in the
world. The practitioner may be internally coherent --- they may do what they
say they will do --- but the effect may still fail to materialise. The goal may
be ill-conceived, the method flawed, the timing wrong. Nature's verdict is:
this action was coherent with intent but incoherent with the world.

This is Failure in the most instructive sense. It does not mean the practitioner
is morally blameworthy; it means they have learned something about the mismatch
between their model of the world and the world itself. The failure is a \emph{gift}
from nature --- the information needed to correct the model.

\section{The Judgment of Coherence}

The judgment of coherence is not made by a judge. It is made by the structure
of the world itself, registered by the \physis\ engine and witnessed by the
\praxis\ community. No one pronounces sentence. The sentence is the consequence:
the action either produced an effect or it did not.

This is what distinguishes natural justice from legal justice. Legal justice is
the judgment of men according to rules (\emph{nomos}). Natural justice is the
judgment of the world according to structure (\physis). The practitioner who
seeks legal justice asks: ``Am I following the rules?'' The practitioner who
seeks natural justice asks: ``Did this produce a real effect in the world?''

The \physis\ engine's coherence tracker is a tool for the second kind of
judgment. It does not punish or reward. It simply \emph{records} --- and the
record is itself the judgment. A life lived in high coherence is a life whose
actions have left traces, produced effects, made differences. A life lived in
low coherence is a life that passed through the world without altering it.

\section{The Socratic Paradox, Resolved}

Socrates famously claimed that no one knowingly does wrong --- that wrongdoing
is always the product of ignorance \parencite{plato1997complete}. The claim has
puzzled philosophers for two millennia, for it seems to deny the reality of weakness
of will \parencite{macintyre1981virtue}. How can Socrates
claim that the person who eats the cake knowing it harms them is simply
\emph{ignorant}?

The coherence framework resolves this paradox. The practitioner who acts against
their own declared intent is not ignorant in the sense of lacking information.
They are \emph{incoherent} --- their action does not match their intent. This
incoherence is registered as -1 by the engine, and it is a form of ignorance:
the practitioner does not yet know what they truly intend. Their declared intent
is not their actual intent; the action reveals the true one.

The discipline of the practitioner is the progressive elimination of this gap
--- the process of bringing declared intent and actual action into alignment.
A life of perfect coherence would be a life in which the practitioner always
does what they declare, and what they declare always produces the intended
effect. Such a life is perhaps unattainable, but it is the ideal toward which
the disciplines tend.

\section{Justice as the Health of Action}

Plato compares justice to health. Health is not a possession but a state: the
body functions as it should, each part doing its own work. Justice, likewise,
is not a possession but a state: the soul functions as it should, each faculty
doing its own work.

We compare coherence to health. Coherence is not a score to maximise but a
state to inhabit. The practitioner in a state of high coherence acts with ease,
with power, with effect --- not because they are forcing themselves but because
their action is aligned with their nature and nature's order. Incoherence is
the disease of action: the will that cannot effect, the intent that cannot
manifest, the life that leaves no trace.

The \physis\ engine's coherence index is thus not a gamification of virtue but
a \emph{diagnostic} --- a measure of the health of action in a given life. A
falling coherence index is a warning sign, like a fever. It tells the practitioner:
something in the relation between your intent and your effect is out of order.
The disciplines of Book III are the remedies for this disorder.

\section{The Deeper Challenge: Schopenhauer's Blind Will}

The account of coherence presented in this book has assumed something significant:
that nature has a direction, a tendency toward +1 over --1, an orientation that
the coherence score registers. But what if the deepest truth of nature is not
coherence but a blind, purposeless striving that cares nothing for success or
failure?

Arthur Schopenhauer's metaphysics of the Will is the most powerful statement of
this challenge. For Schopenhauer, the thing-in-itself --- the reality behind the
veil of appearance --- is not a rational order (as Hegel held) nor a teleological
drive (as Aristotle's entelechy) nor a divine creation. It is \emph{Will}:
a ceaseless, objectless striving that manifests in every phenomenon, from the
gravity that pulls a stone to the desire that drives a human being
\parencite{schopenhauer1818will}. The Will has no purpose, no goal, no
coherence metric. It simply pushes, endlessly, without reason and without rest.

The challenge to our thesis is devastating if sustained. If nature is blind Will,
then coherence (+1/0/--1) is an anthropomorphic projection onto a reality that
has no direction. The practitioner who seeks to align their action with nature
is chasing a phantom, for nature has nothing to align with --- only an endless
striving that is indifferent to the difference between success and failure.

But Schopenhauer's insight is not the final word. He is right that the Will has
no \emph{purpose}, but he is wrong that it has no \emph{structure}. Coherence
is not a purpose the Will intends but a description of what the Will does when
it succeeds. Even blind variation combined with selective retention produces the
appearance of teleology: the organism that finds its niche persists, the one
that does not dies out; the action that produces a real effect is registered,
the one that does not leaves no trace. The +1 is not an intention but a
registration: this form persisted, this pattern worked.

The \physis\ engine's coherence scale is precisely this: not a claim about
nature's intentions but a record of nature's selections. The engine does not ask
``What does nature want?'' It asks ``What has nature preserved?'' The answer
is the coherence score --- not as metaphysics but as data.

But the honest objector will not be satisfied, and we must not let ourselves off
easily. The reply just given purchases its safety at a price: it retreats from
metaphysics to data. ``Coherence is only a record of what persisted'' is true,
but it is also \emph{too weak to do the work the thesis demands of it}. For the
thesis does not merely describe what persists; it commands the practitioner to
\emph{align} with it, to treat the coherent as the standard of right action. And
between ``this pattern persisted'' and ``you ought to pursue this pattern'' there
is a gap that no quantity of data can close. The plague persists. The tumour
persists. The addiction persists --- indeed it is the very type of a self-stabilising
pattern that scores, in crude terms, a relentless +1. If coherence is mere
survival-registration, then the Mirror Trap itself, which is nothing if not a
durable, self-reproducing structure, is maximally coherent. The descriptive
reading of coherence does not merely fail to refute Schopenhauer; it hands him
the field, for it concedes that nature selects without selecting \emph{for the
good}.

We meet the objection by refusing its framing. Schopenhauer's error is not that
he sees the Will as blind but that he treats the Will as \emph{one} ---
a single, undifferentiated striving of which every phenomenon is an equal
expression. This is the move that makes the gravity of the stone and the love of
the human ``the same'' Will at different grades of objectification. But nature,
as we have insisted since Book II, is not one striving; it is a \emph{hierarchy}
of strivings, each with its own form, its own entelechy, its own measure of
what counts as achieving itself. The tumour's persistence is coherence
\emph{at the level of the tumour} and incoherence at the level of the organism
it kills. The addiction's stability is coherence in the reflex arc and the
collapse of the person in whom that arc runs. There is no view from nowhere ---
no flat, level field of bare survival --- from which to read off a single
coherence score. Every +1 is indexed to a form, and forms are nested. This is
why the engine does not score actions in the abstract but always relative to a
declared goal and the holarchy that contains it: coherence is never ``did this
persist?'' but ``did this persist \emph{as the kind of thing it was striving to
be}?''

This nesting is what converts Schopenhauer's blind Will back into something the
practitioner can align with. The practitioner does not align with the Will in
general --- there is no such thing to align with, and Schopenhauer is right that
the attempt would be chasing a phantom. They align with their \emph{own}
entelechy: the specific form their striving is already, implicitly, at work to
become. Schopenhauer saw the striving and missed the form; he universalised the
push and ignored that every push is the push \emph{of some particular being
toward its own actuality}. Restore the form, and the indifference dissolves ---
not because nature now ``cares,'' but because the practitioner is no longer
asking nature at large for a verdict. They are asking whether \emph{this} action
realised \emph{this} nature. That question has a determinate answer, and it is
the coherence score.

\subsection{The Is and the Ought}

We have not yet paid the most expensive debt. Even granting that coherence is
indexed to a being's own form, the sharpest reader will press the classical
objection: this whole edifice slides, without warrant, from \emph{is} to
\emph{ought}. ``This action realised your nature'' is a factual claim; ``you
ought to perform actions that realise your nature'' is a normative one, and no
chain of facts entails a value. To say that the engine's +1 is \emph{justice}
--- and we have said exactly this, calling coherence ``nature's justice'' --- is
to commit the naturalistic fallacy in plain sight: to read a standard of the
good off a description of the world.

We do not deny the gap; we deny that the thesis leaps it illegitimately. The
``ought'' of coherence is not a categorical imperative derived from facts. It is
a \emph{hypothetical} imperative, and its antecedent is supplied not by nature
but by the practitioner. The structure is: \emph{if} you will to be the being
you are --- if you take up your own entelechy as your end --- \emph{then}
coherence is the measure of your success, and incoherence the measure of your
failure. Nature does not command you to live. It offers no premise that obligates
you to anything. The premise is the practitioner's own act of self-appropriation:
the declaration of a goal, which is nothing other than the willing of a form.
The is/ought gap is real, and it is crossed exactly once, by the practitioner,
in the founding act of declaring what they intend to become. Everything
downstream of that act --- every coherence score, every verdict of justice --- is
analytic to it. Justice is not smuggled in from nature; it is unpacked from the
practitioner's own commitment.

This is why the engine cannot, and does not, tell the practitioner what to want.
It is a tribunal with no legislative power. It registers fidelity to a law the
practitioner has given themselves. Plato could ground justice in the Form of the
Good because he held that the Good was seen, not chosen --- that the ascent ended
in a vision that compelled assent. We make no such claim. Our Good is not seen
but \emph{enacted}: it comes into being in the act of commitment and is sustained
by coherence with that commitment. This is a weaker foundation than Plato's, and
we accept the weakness as the price of honesty. The practitioner who declares no
goal is subject to no justice and no failure; they have simply not entered the
game. But the one who does declare --- who says ``this is what I will be'' --- has,
in that very act, made the is/ought gap their own to cross, and the coherence
score is the faithful record of whether they have crossed it.

Schopenhauer's own response to the blind Will --- aesthetic contemplation that
temporarily quiets the Will, and asceticism that denies it entirely --- is one
possible path. But it is the path of the mystic, not the practitioner. The
thesis of this work offers another: not denial of the Will but alignment with
its most coherent expressions. The practitioner-king does not deny the Will;
they ride it as a surfer rides a wave, moving with its structure rather than
resisting it. The wave has no purpose, but it has a shape; the practitioner
who learns the shape can act within it, through it, as it.

\section{On the Nature in Things}

There is a closer ancestor to this whole project than Plato, and naming him will
sharpen what coherence is. Spinoza opens his unfinished \emph{Treatise on the
Emendation of the Intellect} with the resolve to find a method for healing the
mind --- for leading it from inadequate ideas, confused and passively suffered, to
adequate ideas grasped through their causes, and so raising its power of acting
until it rests in the intellectual love of Nature \parencite{spinoza1985emendation}.
The \emph{Ethics} is that method carried out: a geometric ascent from bondage to
the passive affects toward freedom, which is necessity understood
\parencite[V.P42]{spinoza1994ethics}. Praxis is the same project turned to a
different object. Where Spinoza offers an \emph{emendatio intellectus}, a mending
of the intellect, the disciplines of Book III offer an \emph{emendatio actionis},
a mending of action: the same passage from the inadequate to the adequate, from
bondage to freedom, from the passive to the active --- but tracked in deeds and
their effects rather than in ideas and their causes. Coherence is the power of
acting made visible; the rising coherence index is Spinoza's increase of
\emph{conatus} \parencite[III.P6--P7]{spinoza1994ethics}, the striving by which
each thing perseveres in its being, which is nothing other than the entelechy of
Book II seen from within.

But Praxis inverts Spinoza on the very point that defines him. Spinoza is an
intellectualist: emend the idea and right life follows, for the affects are
themselves confused ideas, cured by adequate knowledge, and the mind once purified
is its own judge --- an idea is adequate when it is clear, distinct, and grasped
through its cause, certified by reason alone. Praxis reverses the order of
treatment. It emends through the \emph{deed} and the world's registered response,
not through the clarified idea; its tribunal is not the mind's certainty but
nature's effect (the distinction of internal from external coherence drawn earlier
in this book); and its method is not solitary, geometric, and complete in one
ascent but social, iterative, and never finished --- the witness rather than the
meditation, the dream loop rather than the deduction. Spinoza emends the intellect
and trusts action to follow. Praxis emends action and lets the intellect be
corrected by the world. The one works from the inside out; the other from the
outside in.

And yet the two are not finally opposed, because Spinoza himself supplies the
warrant for their unity. \emph{Ordo et connexio idearum idem est ac ordo et
connexio rerum} --- the order and connection of ideas is the same as the order and
connection of things \parencite[II.P7]{spinoza1994ethics}. Idea and thing are one
substance grasped under two attributes; to mend the order of ideas and to mend the
order of things are therefore one emendation seen from two sides. This is why the
mending of action and the mending of intellect cannot, in the end, come apart: the
nature within things and the nature within thought are the same nature. And here
the engine does something Spinoza could only assert. It builds the parallelism. In
the shared embedding space, a declared intention --- an idea --- and a registered
effect --- a thing --- are projected into one and the same geometry, and coherence
is simply the distance between them. The \physis\ engine is a literal, running
\emph{ordo idearum} that is also an \emph{ordo rerum}: the order of intentions and
the order of effects laid down in a single space so that their agreement can be
measured. What Spinoza proved \emph{more geometrico} the engine computes
\emph{more vectoriali}. The nature in things and the nature in thought meet, at
last, in the same coordinates --- which is only to say, again and from the deepest
ground, that \praxis\ is the tool of \physis.
